% ==================================================================
% §15.5.4 REBUILT — Phase 1: Excluded and stress-test probes
%
% Part 4 of the §15.5 rebuild under the parallel-creation strategy.
% Insert IMMEDIATELY AFTER §15.5.3 block.
%
% GPT round 9 final version (5 corrections applied):
%   (1) Category (ii) closing paragraph softened: "enlarging the
%       phenomenological description beyond the present ε-only layer"
%       (not "incompatibility is structural in absolute sense").
%   (2) BAO paragraph rephrased: methodology-limited, without
%       "most naturally read as an artefact" (too confident).
%   (3) A_lens paragraph rephrased without κ_A ∈ [2,4] numerics.
%   (4) S_8 / σ_8 paragraphs: removed "independent of detailed proxy
%       kernel choice within bracketed range"; rewritten in GPT's
%       uniform-ε-mapping language.
%   (5) SN Ia paragraph made neutral (no "placeholder removed"
%       self-commentary).
%
% NO legacy Hubble compensation numbers anywhere.
% NO specific legacy ε values.
% ==================================================================

\subsubsection*{\normalfont\textbf{15.5.4 \quad Excluded and stress-test probes}}
\label{subsec:eps_excluded_rebuilt}

Beyond the retained five-probe set of \S\ref{subsec:eps_retained_rebuilt},
several additional cosmological channels were considered in the present
analysis but are not used in the joint effective band of
\S\ref{subsec:eps_joint_band_rebuilt}. We group them into three
categories whose reasons for exclusion are structurally different.

\paragraph{Category (i) --- methodology-limited channels.}
These probes address the same effective $\varepsilon$-sector as the
retained probes and could in principle re-enter the joint effective
band at a later stage, but the present extraction pipelines are not
yet explicit enough for headline use. The limitation is therefore
methodological rather than structural.

\begin{itemize}[leftmargin=1.4em,itemsep=2pt,topsep=4pt]
\item \textbf{BAO (DESI 2024 DR1).} The present implementation uses a
simplified shape-only treatment with independently fixed $\Omega_m$,
diagonal covariance, and reduced-grid profiling of $H_0$ and $r_d$.
Under this simplified pipeline the raw $1\sigma$ interval straddles
$\varepsilon = 0$ and has a slightly negative central value. In the
present paper this is treated as a methodology-limited result rather
than as evidence against the $\varepsilon$-sector itself. Headline
inclusion would require a full DESI likelihood with coupled
$(\Omega_m, r_d, \varepsilon)$ fitting and the full covariance
structure. The probe is therefore retained in the reproducibility
pipeline, but excluded from the joint effective band.

\item \textbf{$A_{\rm lens}$ CMB lensing.} The present treatment uses
a linearised proxy of the form
$A_{\rm lens} \approx 1 + \kappa_A\,\varepsilon$ rather than a full
Boltzmann-level calculation. At this level the observational error bar
is not the only uncertainty: the proxy mapping itself carries a
substantial theoretical systematic. The resulting interval is
therefore not robust enough for headline use. Headline inclusion would
require a full \textsc{camb}/\textsc{class}-level analysis. The probe
is retained in the reproducibility pipeline, but excluded from the
joint effective band.
\end{itemize}

\paragraph{Category (ii) --- channels outside the present $\varepsilon$-sector.}
These are distinct in kind from category~(i). In the present
uniform-$\varepsilon$ diagnostic layer their raw $1\sigma$ support
lies entirely at $\varepsilon < 0$. Under the physical prior
$\varepsilon \geq 0$ adopted in \S\ref{subsec:eps_def_rebuilt},
including them as $\varepsilon$-band constraints would therefore be a
category error: the tension they report is not encoded in the present
$\varepsilon$-sector as parameterised here.

\begin{itemize}[leftmargin=1.4em,itemsep=2pt,topsep=4pt]
\item \textbf{$S_8$ weak lensing (KiDS-1000).} In the present
uniform-$\varepsilon$ mapping, the raw $1\sigma$ interval lies
entirely at $\varepsilon < 0$. Within the $\varepsilon$-sector adopted
here this does not act as an $\varepsilon$-band constraint but instead
signals the separate $S_8$ tension, whose resolution is not encoded in
a positive-drift $\varepsilon$ deformation. The probe is therefore
classified as outside the present $\varepsilon$-sector rather than as
a retained $\varepsilon$-channel.

\item \textbf{$\sigma_8$ cluster counts.} The same structural pattern
appears in the present simplified mapping: the raw $1\sigma$ support
lies entirely in $\varepsilon < 0$. As for $S_8$ weak lensing, this
identifies a tension outside the present $\varepsilon$-sector rather
than a retained $\varepsilon$-band constraint. The probe is therefore
excluded on structural, not merely methodological, grounds.
\end{itemize}

Unlike category~(i), the issue here is not merely an incomplete
extraction pipeline. Within the present uniform-$\varepsilon$
diagnostic layer, refining the $\varepsilon$-fit alone would not
convert these probes into retained $\varepsilon$-band constraints:
their raw support points outside the physical $\varepsilon$-sector
adopted in \S\ref{subsec:eps_def_rebuilt}. Any successful treatment
would therefore require enlarging the phenomenological description
beyond the present $\varepsilon$-only layer, rather than simply
improving the same extraction.

\paragraph{Category (iii) --- non-independent or placeholder channels.}

\begin{itemize}[leftmargin=1.4em,itemsep=2pt,topsep=4pt]
\item \textbf{Age $t_0$.} Under self-consistent treatment with
$H_0(\varepsilon)$ inherited from the Primary Hubble channel, the age
test carries no independent $\varepsilon$-information. It is therefore
treated as a downstream self-consistency note rather than as a
retained $\varepsilon$-probe.

\item \textbf{SN~Ia.} No retained SN~Ia $\varepsilon$-channel is
included at the present stage. A genuine SN~Ia probe would require
direct fitting to an observational supernova dataset with the full
calibration and host-systematics chain. Such a pipeline is not yet
implemented in the present analysis, so SN~Ia is not used in the
joint effective band.
\end{itemize}

\paragraph{Cross-category principle.}
Categories (i) and (ii) must be kept distinct. Category~(i) contains
channels that could re-enter the retained set after a sufficiently
upgraded extraction pipeline. Category~(ii) contains channels that,
within the present uniform-$\varepsilon$ sector itself, do not
function as $\varepsilon$-band constraints at all. Keeping this
distinction explicit prevents low-$z$ structure tensions from being
misrepresented as direct $\varepsilon$-band measurements.

\paragraph{Bridge.}
A short comparison of the present ECT effective-$\varepsilon$ result
with alternative approaches to the same observational landscape is
given in \S15.5.5.

% ==================================================================
% END §15.5.4 (rebuilt). §15.5.5--15.5.6 will be inserted here in
% subsequent steps. DO NOT REMOVE existing §15.5 below.
% ==================================================================
